\subsection{الإجابة عن أسئلة البحث}
\label{sec:research-answers}

تتيح الأدلة المعروضة العودة إلى الأسئلة الثلاثة في القسم~\ref{sec:research-questions}. وتربط الإجابات بين الطريقة والتحقق من المُستقبِل ومقارنات الجودة، دون التعامل مع عيناتها المختلفة بوصفها تجربة واحدة.

\subsubsection{السؤال الأول: كيف تحمل الاختيارات البتات؟}
تربط الطريقة البتات بفهارس القائمتين المرتبتين لمواضع الرد وموضوعات التعليق. ويعبّر نموذج الكتابة عن الاختيارين بتعليق عادي عبر \EN{API}، دون الحاجة إلى احتمالات الرموز. وتتيح النصوص المصدرية وقواعد التحضير والترتيب وإعدادات النموذج والحد الزمني للتعليقات تفسير الفهارس نفسها لدى الطرفين. وتحتاج الرسائل المضغوطة كذلك إلى القاموس وتمثيل الحمولة الموافقين. ويجمع نمط التعليقات المتعددة الرسالة الأطول من تعليقاتها الحاملة وفق الترتيب المتفق عليه. وتبين اختبارات الوحدة المشتركة والمثال المسجّل صحة تحويل البتات وإعادة بناء الرسالة بدقة، بينما يشرح القسمان~\ref{sec:receiver-operation} و\ref{sec:reconstruction-validation} مدخلات المُستقبِل والتحقق من تطابق القوائم.

\subsubsection{السؤال الثاني: ما مدى نجاح التوليد والاستخراج؟}
يولّد التنفيذ تعليقات مقبولة ويستخرج رسائل كاملة، ويتعلق نجاحه بالحمولة وشروط الاستقبال. بلغ استخراج مقتطفات المصادر \NUM{\OctExact/\OctAttempts} مع سياق مشترك محضّر. واستخرج اختبار التطوير بحمولات مستقلة \NUM{\DevExact/\DevAttempted} رسائل جرت محاولتها قبل التوقف، مع بقاء \NUM{\DevUnattempted} دون محاولة. وعند جلب سياق الأرشيف بصورة مستقلة، استخرج نمط الكتابة الذي يركز على الموضوع المختار بميزانية سبع محاولات \NUM{\FreshDevSevenExact/\FreshAssigned} رسائل من عينة التطوير و\NUM{\FreshHeldSevenExact/\FreshAssigned} من العينة غير المستخدمة سابقًا في التوليد. هذه عينات منفصلة. وتتيح السجلات مع اختبارات المكونات التمييز بين عمليات النقل والإخفاقات اللغوية؛ فقد منع الخطأ في تحديد الموضوع استخراج رسائل اكتمل توليدها، بينما ترك استنفاد محاولات الكتابة رسائل أخرى غير مكتملة. لذلك يظل قبول المُرسِل واستخراج المُستقبِل نتيجتين منفصلتين. ويوضح القسمان~\ref{sec:october-recovery} و\ref{sec:fresh-archive} شروطهما وأسس عد المحاولات.

\subsubsection{السؤال الثالث: كيف تقارن التعليقات بتعليقات الطريقة المرجعية؟}
تختلف إجابة المحكّم المستقل باختلاف المعيار. ففي عينة المنشورات الـ\NUM{244} المختارة، حققت تعليقات طريقتنا ملاءمة أعلى للأسلوب، واختيرت بوصفها النص الأضعف بمعدل أقل، بينما كان تمييز تعليق \EN{ZLG}~\cite{lin2024zero} بوصفه التعليق المُضاف أصعب. وكانت قيمتا \EN{AUROC} للاشتباه متقاربتين. لذلك تدعم المقارنة نتائج جودة محددة، ولا تمنح ترتيبًا واحدًا للطريقتين عبر جميع المهام. وتحدد العينة المختارة وأحكام النموذج نطاق نتائج الجودة هذه (الجدول~\ref{tab:llm-judge}). أما الاستخراج المساعد الموضح في القسم~\ref{sec:baseline-choice}، فيحد بصورة منفصلة من الاستنتاج بشأن استخراج الرسالة بصورة مستقلة. ولا تثبت مهام تقييم الجودة تفوقًا عامًا في الأمان.
